\chapter{Annexes}
\label{ch:annexes}

\section{Glossaire rapide}

\begin{description}[style=nextline, leftmargin=2.8cm]
  \item[Session] Créneau planifié dans l'agenda, rattaché à un cours et un groupe.
  \item[Plan de charge] Vue d'ensemble des écoles et KPIs annuels.
  \item[Non facturé TTC] Montant des sessions sans facture, TVA incluse (affichage liste écoles).
  \item[Rapprochement] Lien entre une ligne bancaire et une facture, dépense ou note de frais.
  \item[Contexte métier] Profil terminologique education / consulting / medical.
\end{description}

\section{Régénérer les captures d'écran}
\label{ch:annexes-captures}

Les images sont stockées dans \texttt{screenshots/}. Pour les mettre à jour après une évolution de l'interface~:

\begin{enumerate}[leftmargin=*]
  \item Démarrer le serveur local~: \texttt{php artisan serve} à la racine du projet.
  \item S'assurer qu'un compte de démonstration existe (seed ou compte admin).
  \item Exécuter depuis ce dossier~:
\begin{verbatim}
cd documentation/manuel-utilisateur
npm install puppeteer-core   # une seule fois
node scripts/capture-screenshots.mjs
\end{verbatim}
  \item Variables optionnelles~: \texttt{VRP\_SCREENSHOT\_BASE}, \texttt{VRP\_SCREENSHOT\_EMAIL}, \texttt{VRP\_SCREENSHOT\_PASSWORD}.
  \item Recompiler le PDF~: \texttt{make}.
\end{enumerate}

\section{Compiler le PDF}

Prérequis~: \textbf{XeLaTeX} et polices TeX Gyre (distribution TeX Live ou MacTeX).

\begin{verbatim}
cd documentation/manuel-utilisateur
make          # produit manuel-utilisateur-vrp.pdf
make clean    # supprime les fichiers auxiliaires
\end{verbatim}

\section{Documentation technique}

Pour les détails d'implémentation, consultez le dossier \texttt{documentation/fr/} du dépôt~:

\begin{itemize}[leftmargin=*]
  \item \texttt{v2-interface-utilisateur.md}
  \item \texttt{v2-navigation-modules.md}
  \item \texttt{parcours-creation-ecole.md}
  \item \texttt{v2-facturation-par-ecole.md}
  \item \texttt{v2-tresorerie-rapprochement-bancaire.md}
\end{itemize}

\section{Dépannage courant}

\begin{table}[htbp]
  \centering
  \caption{Problèmes fréquents}
  \begin{tabular}{@{}lp{7.5cm}@{}}
    \toprule
    \textbf{Symptôme} & \textbf{Piste de résolution} \\
    \midrule
    Bouton Planifier inactif & Vérifier le mode Édition et l'absence de formulaire imbriqué. \\
    Montants incohérents & Distinction HT (section facturation) vs TTC (liste écoles, trésorerie planifiée). \\
    Import bancaire refusé & Fichier \texttt{.xlsx} avec colonnes attendues, $<$ 10\,Mo. \\
    Pas de suggestion rapprochement & Facture déjà rapprochée ou montant hors tolérance 0{,}02\,€. \\
    \bottomrule
  \end{tabular}
\end{table}

\section{Index des captures}

\begin{table}[htbp]
  \centering
  \begin{tabular}{@{}ll@{}}
    \toprule
    \textbf{Fichier} & \textbf{Écran} \\
    \midrule
    \texttt{01-landing.png} & Page d'accueil publique \\
    \texttt{02-login.png} & Connexion \\
    \texttt{03-demande-acces.png} & Demande de compte \\
    \texttt{04-home.png} & Liste des écoles / plan de charge \\
    \texttt{05-planning.png} & Agenda planning \\
    \texttt{06-treasury.png} & Synthèse trésorerie \\
    \texttt{07-programs.png} & Programmes \\
    \texttt{08-groups.png} & Groupes \\
    \texttt{09-school-billing.png} & Facturation fiche école \\
    \bottomrule
  \end{tabular}
\end{table}
